% Источники: 03_Карточки_доменов/Сессии_задач.md
% Генератор: docs/build_pdf.py
% Шаблон печатного представления. Содержимое подставляет docs/build_pdf.py.
\documentclass[12pt,a4paper]{article}
\usepackage{iftex}
\ifPDFTeX
  \usepackage[utf8]{inputenc}
  \usepackage[T2A]{fontenc}
  \usepackage{newtxtext,newtxmath}
\else
  \usepackage{fontspec}
  \IfFontExistsTF{Times New Roman}
    {\setmainfont{Times New Roman}}
    {\setmainfont{Libertinus Serif}}
  \IfFontExistsTF{Arial}{\setsansfont{Arial}}{\setsansfont{Latin Modern Sans}}
  \IfFontExistsTF{Menlo}{\setmonofont{Menlo}[Scale=MatchLowercase]}
    {\setmonofont{Latin Modern Mono}[Scale=MatchLowercase]}
\fi
\usepackage[russian]{babel}
\usepackage{geometry}
\geometry{left=25mm,right=15mm,top=20mm,bottom=20mm}
\usepackage{setspace}
\onehalfspacing
\usepackage{indentfirst}
\setlength{\parindent}{1.25cm}
\setlength{\parskip}{2pt}
\setlength{\emergencystretch}{3em}
\usepackage{graphicx,booktabs,longtable,array,tabularx,calc,etoolbox}
\usepackage{enumitem}
\setlist[itemize]{leftmargin=1.25cm,itemsep=2pt,topsep=4pt}
\setlist[enumerate]{leftmargin=1.25cm,itemsep=2pt,topsep=4pt}
\newcommand{\tightlist}{\setlength{\itemsep}{0pt}\setlength{\parskip}{0pt}}
\AtBeginEnvironment{longtable}{\small\singlespacing}
\setlength{\LTpre}{8pt}
\setlength{\LTpost}{8pt}
\setlength{\tabcolsep}{5pt}
\setlength{\LTleft}{0pt}
\setlength{\LTright}{0pt}
\usepackage{xurl}
\usepackage[unicode,colorlinks=true,linkcolor=black,urlcolor=blue]{hyperref}
\usepackage{bookmark}
\hypersetup{pdftitle={Карточка домена: сессии задач},pdfauthor={Анохин Г.А.},pdfsubject={Криптотренажёр mephi-12}}
\urlstyle{same}
\pagestyle{plain}
\setcounter{secnumdepth}{3}
\setcounter{tocdepth}{2}

\begin{document}
\hypersetup{pageanchor=false}
\begin{titlepage}
  \thispagestyle{empty}
  \centering
  {\large НИЯУ МИФИ\par}
  \vspace*{3.5cm}
  {\LARGE\bfseries Карточка домена: сессии задач\par}
  \vspace{0.5cm}
  {\large Карточка домена\par}
  \vspace{0.6cm}
  {\large Платформа для изучения и практики курса по криптографии\par}
  {\large Криптотренажёр mephi-12\par}
  \vspace{0.7cm}
  {\large Версия 0.4\par}
  {\large Редакция от 20.11.2026\par}
  \vspace{1cm}
  \begin{tabularx}{\textwidth}{@{}p{3.4cm}Xp{4cm}@{}}
    \toprule
    Роль & ФИО & Подпись / Дата \\
    \midrule
    Разработал & Анохин Г.А. & \\
    \addlinespace[6pt]
    Согласовано & & \\
    \addlinespace[6pt]
    Утверждено & & \\
    \bottomrule
  \end{tabularx}
  \vfill
  {\large Москва, 2026\par}
\end{titlepage}
\hypersetup{pageanchor=true}

\section*{Лист регистрации изменений}
\addcontentsline{toc}{section}{Лист регистрации изменений}
\begingroup\small\singlespacing
\noindent
\begin{tabularx}{\textwidth}{@{}p{1.2cm}p{2cm}p{2.5cm}Xp{1.8cm}@{}}
\toprule
Версия & Дата & Автор & Содержание изменения & Утвердил \\
\midrule
0.4 & 20.11.2026 & Анохин Г.А. & Описаны автоматическое закрытие мероприятия, срок личной сессии и повторное прохождение по настройке преподавателя. & \\
\bottomrule
\end{tabularx}
\endgroup
\clearpage
\begingroup\small\singlespacing
\tableofcontents
\endgroup
\clearpage

\hypertarget{карточка-домена-сессии-задач}{}

Версия 0.4 от 20.11.2026. Владелец домена --- Анохин Г.А.

\section{Назначение и границы}\label{ux43dux430ux437ux43dux430ux447ux435ux43dux438ux435-ux438-ux433ux440ux430ux43dux438ux446ux44b}

\hypertarget{назначение-и-границы}{}

Сессия задач --- запись одного прохождения работы одним студентом. Она хранит выданные ему условия, время начала, ответы и результаты. Владелец сессии --- студент, для которого она создана. Эта карточка определяет, когда он может сдать ответ и как узнать исход проверки.

Для контрольной работы преподаватель выбирает задания и их порядок, а также указывает, разрешено ли повторное прохождение. Каждое место задания в списке --- \textbf{позиция}. Пока работа находится в черновике, она скрыта от студентов. Преподаватель публикует работу для одной из групп, которые назначил ему администратор, и задаёт дату и время открытия и закрытия. Состав закрепляется при публикации. Студент видит назначение в расписании группы, подтверждённой через его учётную запись МИФИ. Когда наступает время открытия, студент может начать работу и получить по каждой позиции отдельную задачу. Для самостоятельной практики источником набора служит \textbf{шаблон} --- именованный список типов заданий, доступный группе. Правила подготовки определены в {\href{../../02\_Техническое\_задание/Техническое\_задание.pdf\#ts-req-func-event-create}{\texttt{TS.\allowbreak{}REQ.\allowbreak{}FUNC.\allowbreak{}EVENT.\allowbreak{}CREATE}}}, {\href{../../02\_Техническое\_задание/Техническое\_задание.pdf\#ts-req-func-event-publish}{\texttt{TS.\allowbreak{}REQ.\allowbreak{}FUNC.\allowbreak{}EVENT.\allowbreak{}PUBLISH}}} и {\href{../../02\_Техническое\_задание/Техническое\_задание.pdf\#ts-req-func-config-view}{\texttt{TS.\allowbreak{}REQ.\allowbreak{}FUNC.\allowbreak{}CONFIG.\allowbreak{}VIEW}}}.

\textbf{Тип задания} определяет правила выдачи условия и проверки решения. Их выполняет соответствующий модуль программы. В ранцевых задачах и упражнениях по схеме Эль-Гамаля модуль рассчитывает учебные параметры. Для текстового задания (\texttt{CUSTOM\_\allowbreak{}TEXT}) преподаватель заранее вводит одну или несколько пар «условие --- правильный ответ». Этот ответ называется эталоном. Для каждой позиции с таким заданием система случайно выбирает пару при старте сессии. Студент видит условие выбранного варианта при каждом открытии своей задачи. Разным студентам может выпасть один вариант. Сравнение ответа с эталоном описано в {\href{../../02\_Техническое\_задание/Техническое\_задание.pdf\#ts-req-func-task-assign}{\texttt{TS.\allowbreak{}REQ.\allowbreak{}FUNC.\allowbreak{}TASK.\allowbreak{}ASSIGN}}}.

\textbf{Вердикт} --- сохранённый результат содержательной проверки ответа одной задачи. Попытка расходуется при записи вердикта «верно» или «неверно». Общее состояние сессии учитывает все выданные задачи. Таблица {\href{../methodology/Глоссарий.pdf\#ответ-попытка-и-состояние}{состояний задачи и сессии}} объясняет различия. Программная роль студента обозначается \texttt{USER}.

\section{Сценарии}\label{ux441ux446ux435ux43dux430ux440ux438ux438}

\hypertarget{сценарии}{}

{\def\LTcaptype{none} % do not increment counter
\begin{longtable}[]{@{}
  >{\raggedright\arraybackslash}p{(\linewidth - 2\tabcolsep) * \real{0.3200}}
  >{\raggedright\arraybackslash}p{(\linewidth - 2\tabcolsep) * \real{0.6800}}@{}}
\toprule\noalign{}
\begin{minipage}[b]{\linewidth}\raggedright
Операция владельца с ролью \texttt{USER}
\end{minipage} & \begin{minipage}[b]{\linewidth}\raggedright
Что получает студент
\end{minipage} \\
\midrule\noalign{}
\endhead
\bottomrule\noalign{}
\endlastfoot
Начать контрольное мероприятие по идентификатору \texttt{assessmentId} из расписания & Личный набор по опубликованному составу для своей группы. Начать работу можно с момента открытия и строго до закрытия. Для просмотра прежней сессии используется её собственный идентификатор \\
Начать самостоятельную практику по имени шаблона & Полный набор задач в порядке, заданном выбранным шаблоном для группы студента. Если у студента уже есть незавершённая сессия (\texttt{NEW}), отказ содержит её идентификатор. Он позволяет открыть набор, когда ответ на создание потерян \\
Отправить решение по идентификатору задачи & При успешной обработке --- сохранённый вердикт по конкретной задаче. Проверяющий использует её выданное условие и тип. На задачу допускается одна предметная попытка \\
Прочитать собственную сессию по идентификатору & Текущее состояние набора и каждой задачи, включая сохранённые вердикты. Чтение доступно и после общего срока. Чужая и отсутствующая сессии дают одинаковый отказ \\
\end{longtable}
}

Если при старте контрольной потеряна связь, студент снова открывает расписание. Карточка содержит идентификатор его активной сессии по этой работе, а после её завершения --- последней собственной сессии. По этому идентификатору можно открыть прежние условия и результаты. После завершения студент может начать новое прохождение, только если преподаватель разрешил повтор и окно ещё открыто. Новая сессия получает свой идентификатор и варианты. Закрытие мероприятия прекращает выдачу новых сессий. Студент может по-прежнему читать свои прежние результаты.

Практика и контроль различаются по записи сессии. В самостоятельной практике поле \texttt{assessmentId} равно \texttt{null} (связи с мероприятием нет). Результат такой работы остаётся личным результатом студента и не засчитывается за мероприятие, даже если оба набора имеют одинаковый состав.

При получении полного запроса сервер отмечает время его поступления и устанавливает принадлежность задачи подтверждённому отправителю. Для собственной задачи запрос регистрируется с исходным временем до ожидания в очереди. Попытка обратиться к чужой или отсутствующей задаче завершается отказом. Такое обращение не открывает операцию и не задерживает завершение чужой задачи.

Своевременная регистрация сохраняет возможность пройти проверки допуска и формы, пока действует предел обработки. Право начать предметную проверку получает один допустимый ответ на задачу. Если зарегистрированы два запроса, их тела сохраняют свою принадлежность. Второй запрос не может заменить ответ, уже находящийся на предметной проверке.

\section{Срок и попытка}\label{ux441ux440ux43eux43a-ux438-ux43fux43eux43fux44bux442ux43aux430}

\hypertarget{срок-и-попытка}{}

Время создания набора на сервере обозначается \texttt{t0}, а граница приёма новых ответов --- \texttt{D}. Все задачи одной сессии имеют общий срок. Для самостоятельной практики \texttt{D\ =\ t0\ +\ 20\ минут}. Для контрольной сервер выбирает более ранний момент: истечение 20 минут после старта или назначенное преподавателем время закрытия \texttt{closesAt}. Например, при старте в 10:25 и закрытии в 10:30 полный ответ должен поступить до 10:30. Ровно в 10:30 сдавать уже поздно.

Время получения полного запроса обозначается \texttt{t\_in}. Предел записи вердикта по нему --- \texttt{L =\allowbreak{} t\_\allowbreak{}in + 60 секунд}. В этот интервал входят очередь, вычисление, возможное восстановление и запись результата. Сервер может закончить проверку ответа, поступившего до закрытия, и после закрытия, если успевает записать вердикт строго до \texttt{L}. Планировщик раз в 60 секунд ищет просроченные задачи. Расписание этих обходов не меняет \texttt{D} и \texttt{L}.

{\def\LTcaptype{none} % do not increment counter
\begin{longtable}[]{@{}
  >{\raggedright\arraybackslash}p{(\linewidth - 2\tabcolsep) * \real{0.3200}}
  >{\raggedright\arraybackslash}p{(\linewidth - 2\tabcolsep) * \real{0.6800}}@{}}
\toprule\noalign{}
\begin{minipage}[b]{\linewidth}\raggedright
Событие
\end{minipage} & \begin{minipage}[b]{\linewidth}\raggedright
Последствие для задачи
\end{minipage} \\
\midrule\noalign{}
\endhead
\bottomrule\noalign{}
\endlastfoot
Полный запрос получен при \texttt{t\_in\ \textless{}\ D} & Для зарегистрированного ответа владельца можно завершить допуск и проверку. Вердикт разрешено сохранить после \texttt{D}, если фиксация произошла строго до \texttt{L} \\
Полный запрос владельца получен при \texttt{t\_in\ \textgreater{}=\ D} & Новый ответ на собственную задачу отклоняется по сроку, даже если фоновая проверка ещё не закрыла задачу \\
К моменту \texttt{L} вердикт не зафиксирован & Операция завершается техническим отказом. Результат, поступивший в \texttt{L} или позже, не записывается \\
Формат неверен или обработка завершилась техническим отказом без вердикта & Попытка сохраняется. Исправленный или новый запрос допустим только при собственном \texttt{t\_in\ \textless{}\ D} и соблюдении остальных условий \\
Наступил \texttt{D}, у незавершённой задачи нет своевременного ответа владельца в обработке & Задача подлежит завершению по просрочке. Её состояние --- \texttt{FAILED}, предметный результат отсутствует \\
\end{longtable}
}

Попытка расходуется при фиксации вердикта «верно» или «неверно». Перезапуск и повторная отправка не продлевают общий срок \texttt{D} или предел \texttt{L} ранее принятой операции. Новый допустимый запрос получает собственное время поступления \texttt{t\_in}. После перезапуска незавершённая операция продолжается по тем же входным данным до прежнего \texttt{L}. Если вердикт уже зафиксирован, студент читает его в своей сессии. После окончательного технического отказа та же операция уже не может записать вердикт.

\section{Итог сессии}\label{ux438ux442ux43eux433-ux441ux435ux441ux441ux438ux438}

\hypertarget{итог-сессии}{}

Пока хотя бы одна задача не завершена, сессия остаётся \texttt{NEW}. Это относится и к своевременному ответу, который ещё проверяется. Когда последняя задача завершена, сессия получает \texttt{SOLVED}, если все задачи решены верно, либо \texttt{FAILED}, если хотя бы одна решена неверно или просрочена. Ранее сохранённые вердикты остальных задач остаются доступны.

У отдельной задачи состояние \texttt{FAILED} встречается в двух случаях. Неверное решение имеет сохранённый отрицательный вердикт с пояснением. Если истёк срок, а предметного вердикта не было, поле \texttt{result} равно \texttt{null}. Чтение различает эти исходы и не расходует попытку.

\section{Связанные материалы}\label{ux441ux432ux44fux437ux430ux43dux43dux44bux435-ux43cux430ux442ux435ux440ux438ux430ux43bux44b}

\hypertarget{связанные-материалы}{}

{\href{../../04\_Системные\_требования/DOC-05\_Системные\_требования.pdf\#doc-05-требования-к-сессиям-задач}{Системные требования DOC-05}} задают обязательства перед студентом. {\href{../../07\_Программные\_требования/DOC-06.2\_Сессии\_задач\_и\_планировщик.pdf\#doc-062-сессии-задач-и-планировщик}{DOC-06.2}} определяет действия обработчиков и планировщика. {\href{../../07\_Программные\_требования/DOC-06.2\_Сессии\_задач\_и\_планировщик.pdf\#doc-062-контракты-сессий-задач}{Приложение с контрактами}} содержит схемы данных.

В {\href{../../09\_Архитектура\_и\_API/DOC-08\_Справочник\_API.pdf\#doc-08-справочник-api-сессии-задач}{разделе API DOC-08}} приведены запросы и ответы для трёх операций. {\href{../../10\_Верификация\_и\_тестирование/DOC-10\_План\_верификации.pdf\#doc-10-верификация-сессий-задач-и-планировщика}{Сценарии DOC-10}} описывают, как проверить обе временные границы, восстановление и чтение результата.

{\def\LTcaptype{none} % do not increment counter
\begin{longtable}[]{@{}
  >{\raggedright\arraybackslash}p{(\linewidth - 2\tabcolsep) * \real{0.3200}}
  >{\raggedright\arraybackslash}p{(\linewidth - 2\tabcolsep) * \real{0.6800}}@{}}
\toprule\noalign{}
\begin{minipage}[b]{\linewidth}\raggedright
Дата
\end{minipage} & \begin{minipage}[b]{\linewidth}\raggedright
Изменение
\end{minipage} \\
\midrule\noalign{}
\endhead
\bottomrule\noalign{}
\endlastfoot
05.11.2026 & Описаны общий срок, единственная попытка и момент получения ответа \\
06.11.2026 & Установлены предел обработки и правила восстановления \\
20.11.2026 & Описаны начало, повтор и закрытие контрольной, отдельный срок для практики и предел записи результата \\
\end{longtable}
}


\end{document}
